\section{Approach}
\label{sec:approach}

\subsection{Overview}
\label{sec:overview}

\draft{The approach connects a mechanized specification, an agent-proposed property, and a test of an implementation (Figure~\ref{fig:overview}).}
\draft{The property is written in a restricted SpecTec-like language and interpreted in Lean over the translated specification.}
\draft{The intended workflow derives three artifacts from this property: a Lean proposition, a generator of inputs satisfying the premise, and an executable checker for the conclusion.}
\draft{The property and generator are translated to Lean, while the result checker is an external mechanical procedure.}
\draft{The agent writes proofs for the property and generator, which Lean checks; the generator and result checker are then used for property-based testing of the implementation.}

\begin{figure}[t]
  \centering
  % One Lean context contains both mechanically derived and agent-authored code.
% Orange arrows: agent construction. Blue arrows: Lean proof checking only.
% Dashed arrows: mechanical translation or derivation. Solid black arrows: testing data flow.
\begin{tikzpicture}[
  font=\fontsize{9}{11}\selectfont,
  box/.style={draw,rounded corners=2pt,align=center,minimum height=8mm,inner sep=3pt},
  input/.style={box,rounded corners=0pt,thick,fill=white,text width=2.85cm,minimum height=1cm},
  doc/.style={chamfered rectangle,chamfered rectangle corners=north east,
    chamfered rectangle xsep=1.8mm,draw,thick,align=center,minimum height=9mm,inner sep=3pt},
  leandoc/.style={doc,draw=blue!60!black,fill=blue!9},
  parsed/.style={box,draw=black!80,thick,fill=white,minimum width=1.8cm},
  edge/.style={-{Stealth[length=1.8mm]},thick},
  agentarrow/.style={edge,very thick,draw=orange!75!black},
  leanarrow/.style={edge,very thick,draw=blue!60!black},
  scriptarrow/.style={edge,dashed,draw=black!75},
  lbl/.style={font=\fontsize{8}{9.5}\selectfont,align=center,inner sep=2pt,text=black},
  agentlbl/.style={lbl,text=orange!75!black,font=\fontsize{8}{9.5}\selectfont\bfseries},
  leanlbl/.style={lbl,text=blue!60!black,font=\fontsize{8}{9.5}\selectfont\bfseries},
  heading/.style={font=\fontsize{9}{11}\selectfont\bfseries,anchor=south west,text=black}
]
  % The frame denotes the language of the artifacts, not who creates them.
  \begin{scope}[on background layer]
    \draw[draw=black!55,thick,fill=black!2] (-2.3,-2.55) rectangle (1.1,3.75);
    \draw[rounded corners=4pt,draw=blue!40,fill=blue!4] (3.9,-2.55) rectangle (12.75,3.75);
  \end{scope}
  \node[heading] at (-2.1,3.25) {Inputs};
  \node[heading,text=blue!60!black] at (4.05,3.25) {Lean definitions, propositions, and proofs};

  % Source inputs remain plain rectangular boxes.
  \node[input] (sourcespec) at (-0.6,2.5) {SpecTec\\specification};
  \node[input] (sourceproperty) at (-0.6,0.35) {Property in\\SpecTec-like notation};
  \node[input] (implementation) at (-0.6,-1.9) {Implementation};
  \node[parsed] (parsedproperty) at (2.85,0.35) {Parsed\\property};
  \draw[scriptarrow] (sourceproperty.east) -- node[lbl,above] {parse} (parsedproperty.west);

  \node[leandoc,minimum width=3cm] (specification) at (6.0,2.5) {Specification in Lean};
  \node[leandoc,minimum height=1.15cm] (property) at (6.0,0.35)
    {Property in Lean\\[2pt]$\mathit{Precondition}(x)\Rightarrow\mathit{Goal}(x)$};
  \node[leandoc,minimum width=3cm] (generatorprop) at (10.5,0.35)
    {Generator sound and\\complete for\\$\mathit{Precondition}(x)$};
  \draw[scriptarrow] (sourcespec.east) -- node[lbl,above] {mechanically translate} (specification.west);

  % All three outputs come from the parsed source property, not the Lean proposition.
  \draw[scriptarrow] (parsedproperty.east) -- (property.west);
  \draw[scriptarrow] (parsedproperty.north) -- (2.85,1.55)
    -- node[lbl,above,pos=0.72] {from Precondition} (10.5,1.55) -- (generatorprop.north);

  \node[leandoc,minimum width=1.65cm] (propertyproof) at (6.0,-1.9) {Property\\proof};
  \node[leandoc,minimum width=1.65cm] (generatorproof) at (9.4,-1.9) {Generator\\proofs};
  \node[leandoc,minimum width=1.65cm] (generator) at (11.7,-1.9) {Input\\generator};
  \draw[agentarrow] ([xshift=-2.5mm]property.south) -- node[agentlbl,left] {agent\\proves} ([xshift=-2.5mm]propertyproof.north);
  \draw[leanarrow] ([xshift=2.5mm]propertyproof.north) -- node[leanlbl,right] {Lean\\checks} ([xshift=2.5mm]property.south);
  % The agent constructs both code and proofs; only proofs establish the propositions.
  \draw[very thick,draw=orange!75!black] (generatorprop.south)
    -- node[agentlbl,right] {agent\\constructs} (10.5,-1.0);
  \draw[agentarrow] (10.5,-1.0) -| ([xshift=2.5mm]generatorproof.north);
  \draw[agentarrow] (10.5,-1.0) -| (generator.north);
  \draw[leanarrow] ([xshift=-2.5mm]generatorproof.north)
    -- node[leanlbl,left] {Lean\\checks} ([xshift=-2.5mm]generatorproof.north |- generatorprop.south);

  % The checker receives the generated test inputs and the supplied implementation.
  \node[doc,draw=black,fill=white,minimum width=2.6cm] (checker) at (11.7,-4.0)
    {Checker\\for $\mathit{Goal}(x)$};
  \draw[scriptarrow] (parsedproperty.south) -- (2.85,-4.0)
    -- node[lbl,above,pos=0.7] {from Goal} (checker.west);
  \node[box,minimum width=1.7cm] (verdict) at (14.4,-4.0) {Pass / fail};
  \draw[edge] (generator.south) -- node[lbl,right] {sample\\test inputs} (checker.north);
  \draw[edge] (implementation.east) -- (1.45,-1.9) -- (1.45,-5.0) -- (11.7,-5.0) -- (checker.south);
  \draw[edge] (checker.east) -- (verdict.west);

  % Arrow styles describe operations independently of the Lean artifact frame.
  \draw[scriptarrow] (-2.1,-5.8) -- (-1.3,-5.8);
  \node[lbl,anchor=west] at (-1.2,-5.8) {Mechanical translation / derivation (\texttt{spec4pbt})};
  \draw[edge] (7.9,-5.8) -- (8.7,-5.8);
  \node[lbl,anchor=west] at (8.8,-5.8) {Runtime data flow};
\end{tikzpicture}

  \caption{\draft{Overview of the proposed approach. Our SpecTec-to-Lean translator supplies the Lean specification used for proof checking (not shown). From the SpecTec specification, the agent writes a property in SpecTec-like notation. The workflow translates (1) the property and (2) its input generator into Lean (blue), and (3) the result checker into an external mechanical procedure (purple). The agent writes Lean proofs for the property and generator (orange), which Lean checks. For implementation PBT, (2) supplies inputs to the implementation and (3) checks the observed results.}}
  \Description{A SpecTec specification guides an agent-written SpecTec-like property. The property branches into a blue Lean property, a blue Lean input generator, and a purple external result checker. Orange text in the first two boxes marks agent-written proofs checked by Lean. The generator feeds inputs to the implementation, whose observed results are judged by the result checker to produce pass or fail verdicts. The separate SpecTec-to-Lean translation of the specification is not shown.}
  \label{fig:overview}
\end{figure}

\subsection{User Inputs}
\label{sec:inputs}

\draft{The user supplies the specification and implementation, together with code that serializes generated inputs and extracts observations from implementation outputs.}
\draft{This connection determines what the test observes and remains outside the Lean proofs.}
% Clarify which adapters are supplied by the user and validate their correspondence.

\subsection{From SpecTec to Lean}
\label{sec:translation}

\draft{We introduce a translator from the SpecTec constructs used by P4-SpecTec into Lean definitions and propositions on which the test obligations can be stated.}
\draft{The translator and its output are part of the trusted base: Lean checks proofs about the translated specification, but does not establish that the translation preserves the source specification's meaning.}
\draft{The translation has passed validation tests, described in Section~\ref{sec:validation}.}
% TODO: Specify the supported SpecTec fragment, translation rules, generated Lean declarations, and treatment of unsupported constructs.

\subsection{Property Language and Admissibility}
\label{sec:propose}

\draft{The agent proposes a property in a restricted SpecTec-like language, with a premise describing test inputs and a conclusion describing expected observations.}
\draft{Each property is interpreted as a Lean proposition over the translated specification.}
\draft{The language is intended to restrict properties to forms that support both generation of inputs satisfying the premise and mechanical decision of the conclusion on implementation observations.}
\draft{These requirements concern how properties are expressed and executed; proofs of their correctness are carried out in Lean.}

\draft{The restricted language also provides a common notation for users to inspect agent-proposed properties against the specification.}
\draft{Keeping input conditions and expected observations explicit is intended to make this review easier for readers familiar with the specification.}
\draft{The property interface is therefore organized around statements with a defined testing interpretation, while Lean provides the general language for their proofs.}
% TODO: Illustrate this review process once the grammar is fixed; improved user comprehension is a design motivation, not an evaluated result.
% TODO: Define the concrete grammar and admissibility conditions once the design is fixed. Separate generated inputs, observed outputs, and auxiliary variables, and specify supported relations, binding and execution modes, quantification, and recursion.
% TODO: Give accepted and rejected examples. SpecTec-like notation alone does not establish that a suitable generator or terminating decision procedure can be derived.

\subsection{Deriving the Test Artifacts}
\label{sec:artifacts}

\draft{For a property admitted by the language, the intended workflow derives the three artifacts in Figure~\ref{fig:overview} from the same source.}
\draft{First, it translates the property into a Lean proposition over the mechanized specification.}
\draft{Second, it uses the premise to construct an executable generator of inputs satisfying that premise.}
\draft{Third, it derives a result checker from the conclusion, to decide whether the implementation's observations meet the stated expectation.}
\draft{The workflow also states the corresponding correctness obligations in Lean, which must be discharged before these artifacts are used for testing.}
% TODO: Define the property language's Lean interpretation and its correspondence with generator and oracle derivation; distinguish proved transformations from trusted ones.
% TODO: Specify the generation procedure, any required relation-specific support, and the treatment of size bounds and failed generation.
% TODO: Define the conclusion-to-oracle translation, its termination requirements, and the treatment of resource exhaustion.

\subsection{Checking the Proof Obligations}
\label{sec:verification}

\draft{The agent supplies proofs of the translated property and the generator's correctness, and Lean checks the submitted terms.}
\draft{The property proof establishes the claimed behavior over the specification.}
\draft{The generator proofs establish soundness, meaning that generated inputs satisfy the premise, and completeness, meaning that every input satisfying it is within the generator's support.}
\draft{A proof attempt may expose a false property or a missing class of generated inputs before implementation tests are run.}
% TODO: Fix the precise completeness claim, including how it depends on generation bounds.
% Specify proof checker policy and example failures after the implementation and study are fixed.

\draft{For the result checker, the required agreement is that its abstract execution machine returns a verdict matching the conclusion on the supplied inputs and observations.}
\draft{We describe the procedure and machine in Lean and require a proof of this agreement before using the checker as the test oracle.}
% TODO: Decide which oracle correctness proofs are shared by the translation and which are required for individual properties.

\subsection{Running Property-Based Tests}
\label{sec:test}

\draft{Once the required proofs have been checked, the generator and result checker form the executable test pipeline.}
\draft{A script samples an input from the generator, serializes it using the user-supplied adapter, and runs the implementation.}
\draft{It then extracts the observations and applies the result checker to obtain a verdict.}
\draft{A failure then points to a disagreement with the specification, subject to the tested translation and the trusted input/output connection.}
% The P4 observation and oracle pipeline remains to be implemented.
